\chapter{Healthcare Professionals}

\section*{Definition and Overview}
Healthcare professionals are the people who provide care across the health system, including doctors, nurses, pharmacists, allied health professionals, and support workers. In Australia, healthcare professionals are regulated by the Australian Health Practitioner Regulation Agency (AHPRA), which ensures they are qualified, registered, and meet professional standards. The team of professionals involved in a person's care varies with their needs: GPs coordinate general care, specialists provide expertise in specific areas, and allied health professionals such as physiotherapists, psychologists, and dietitians support rehabilitation, mental health, and wellbeing. Knowing who provides which services helps people use the health system effectively.

\section*{Epidemiology}
The health workforce is large and growing. Australia has around 40,000 GPs, 30,000 specialists, 400,000 nurses and midwives, and hundreds of thousands of allied health professionals. Around 90\% of Australians see a GP each year, and most people see several health professionals annually. There are workforce shortages in rural and remote areas and in some specialties, which affects access to care. The health workforce is increasingly multidisciplinary, with nurses, pharmacists, and allied health professionals taking on expanded roles in providing care.

\section*{Aetiology and Pathophysiology}
This chapter concerns the roles of health professionals rather than a disease. The main groups include:

\begin{itemize}
    \item \textbf{GPs (general practitioners):} provide primary care, coordinate care, manage chronic conditions, and refer to specialists
    \item \textbf{Specialist doctors:} diagnose and treat conditions in specific areas, such as cardiology, surgery, and psychiatry
    \item \textbf{Nurses and midwives:} provide clinical care, education, and support in hospitals and the community
    \item \textbf{Pharmacists:} dispense medicines and provide medicine advice and health services
    \item \textbf{Allied health professionals:} physiotherapists, occupational therapists, psychologists, dietitians, podiatrists, speech pathologists, and others
    \item \textbf{Aboriginal and Torres Strait Islander health workers:} provide culturally safe care
    \item \textbf{Support workers:} assist with daily living, home care, and personal care
    \item \textbf{Paramedics and emergency workers:} provide pre-hospital emergency care
\end{itemize}

\section*{Clinical Features}
People engage with different health professionals depending on their needs.

\begin{itemize}
    \item Routine check-ups, acute illness, and chronic disease management with a GP
    \item Specialist referral for complex or specific conditions
    \item Hospital care from doctors, nurses, and allied health teams
    \item Medicine supply and advice from pharmacists
    \item Rehabilitation and therapy from allied health professionals
    \item Mental health care from psychologists, psychiatrists, and counsellors
    \item Maternity care from midwives and obstetricians
    \item Community care from nurses and support workers
\end{itemize}

\section*{Differential Diagnosis}
Choosing the right professional matters for effective care.

\begin{itemize}
    \item GP versus specialist: most care starts with a GP, who refers to specialists when needed
    \item Self-care versus professional care: many minor conditions are managed at home or with a pharmacist
    \item Allied health versus medical care: allied health professionals provide therapy and management within their scope
    \item Public versus private care: services differ in cost and waiting times
    \item Registered versus unregistered practitioners: registration with AHPRA or another regulator indicates recognised qualifications
    \item Telehealth versus in-person care: both are valid options
\end{itemize}

\section*{When to Seek Medical Care}
See a GP for new, persistent, or concerning symptoms, for health checks, and to coordinate your care. Your GP can refer you to the right specialist or allied health professional.

\textbf{Call 000 (or local emergency services) for:}
\begin{itemize}
    \item Chest pain, severe breathing difficulty, stroke symptoms, or major trauma
    \item Severe bleeding, collapse, or loss of consciousness
    \item Any sudden, severe symptom that worries you
    \item Emergency departments and paramedics provide urgent care regardless of which professionals you usually see
\end{itemize}

\section*{Diagnosis and Investigations}
Working effectively with health professionals involves preparation.

\begin{itemize}
    \item \textbf{Choosing a regular GP:} continuity of care improves health outcomes
    \item \textbf{Checking registration:} AHPRA's public register confirms a professional is registered
    \item \textbf{Bringing information:} medicines lists, test results, and a summary of concerns
    \item \textbf{Preparing questions:} writing down questions before appointments
    \item \textbf{Understanding referrals:} asking what a referral is for and what will happen
    \item \textbf{Coordinating care:} ensuring your GP knows about all care you receive
    \item \textbf{Seeking a second opinion:} reasonable when a diagnosis or treatment plan is uncertain
\end{itemize}

\section*{Management}
Getting the most from health professionals means being an active partner in your care.

\begin{itemize}
    \item \textbf{Use your GP as the coordinator:} they see the whole picture and coordinate referrals
    \item \textbf{Ask questions:} about your diagnosis, treatment options, and what to expect
    \item \textbf{Share information:} tell every professional about your medicines, allergies, and conditions
    \item \textbf{Follow treatment plans:} take medicines as prescribed and attend follow-up
    \item \textbf{Request referrals:} ask your GP about specialist and allied health options
    \item \textbf{Use care plans:} Medicare-funded plans provide subsidised allied health sessions
    \item \textbf{Provide feedback:} complaints and compliments improve care
    \item \textbf{Respect the team:} different professionals work together for your benefit
\end{itemize}

\section*{Complications}
\begin{itemize}
    \item Fragmented care when information is not shared between professionals
    \item Duplication of tests and medicines when care is not coordinated
    \item Delayed diagnosis when symptoms are not taken seriously
    \item Misunderstandings about roles and referrals
    \item Out-of-pocket costs when the wrong service is chosen
    \item Unregulated practitioners providing unsafe care
    \item Discontinuity when changing GPs or specialists
\end{itemize}

\section*{Prognosis and Outcomes}
People who have a regular GP and an engaged, coordinated care team have better health outcomes, including fewer hospitalisations and better management of chronic conditions. The Australian health system provides access to a wide range of well-trained, regulated professionals, and most people's experiences are positive. Being an informed, prepared, and communicative patient improves the value and safety of the care you receive.

\section*{Prevention and Self-Care}
\begin{itemize}
    \item Establish a relationship with a regular GP
    \item Attend recommended health checks and screening
    \item Prepare for appointments with questions and information
    \item Keep a current list of your medicines and allergies
    \item Ask about bulk billing and costs before appointments
    \item Check that health professionals are registered with AHPRA
    \item Follow treatment and follow-up plans
    \item Seek a second opinion if you are unsure about a diagnosis or plan
    \item Provide feedback to improve services
\end{itemize}

\section*{Sources and Further Reading}
The following reputable sources provide further information for healthcare students and patients seeking reliable, current advice about healthcare professionals.

\begin{itemize}
    \item Australian Health Practitioner Regulation Agency. \textit{Check a practitioner's registration.} https://www.ahpra.gov.au/
    \item Healthdirect Australia. \textit{Health professionals and services.} https://www.healthdirect.gov.au/
    \item Royal Australian College of General Practitioners. \textit{What is a GP?} https://www.racgp.org.au/
    \item Australian Nursing and Midwifery Federation. \textit{Nursing and midwifery roles.} https://www.anmf.org.au/
    \item Services Australia. \textit{Allied health and Medicare.} https://www.servicesaustralia.gov.au/
    \item World Health Organization. \textit{Health workforce.} https://www.who.int/
\end{itemize}
